\documentclass[14pt,aspectratio=169,compress]{beamer}
\usetheme{default}
% something tries to renewcommand these but beamer must break the assumption that they are defined so I am defining them
\newcommand{\labelitemi}{}
\newcommand{\labelitemii}{}
\newcommand{\labelitemiii}{}
\newcommand{\labelitemiv}{}

\usepackage{greg}

\useoutertheme[subsection=false]{miniframes}
\useinnertheme{circles}
\setbeamertemplate{navigation symbols}{}
\setbeamercovered{transparent}
\definecolor{cmu_red}{RGB}{153,0,0} % #990000
\definecolor{cmu_dark_red}{RGB}{100,0,0} % #990000
\definecolor{cmu_dark_grey}{RGB}{70, 70, 70} % #464646
\usecolortheme[named=cmu_red]{structure}
\setbeamercolor{palette secondary}{bg=cmu_dark_grey,fg=white}
\setbeamercolor{palette tertiary}{bg=cmu_dark_red,fg=white}
\setbeamercolor{section in head/foot}{bg=cmu_red,fg=white}

\title[Geometry in HoTT]{Discrete differential geometry in homotopy type theory}
\author{Greg Langmead}
\institute[CMU]{Carnegie Mellon University}
\logo{\includegraphics[width=100pt]{figs/cmu-wordmark-horizontal-r.pdf}}
\date{April 2025}

\begin{document}
\begin{frame}
\titlepage
\end{frame}

\begin{frame}
\tableofcontents
\end{frame}

% \begin{frame}
% \begin{itemize}
% \item 
% \item 
% \end{itemize}
% \end{frame}

\section{Intro}

\begin{frame}{Motivation}
Chern-Weil theory is an algorithm for manifold \( M \), Lie group \( G \)
\begin{enumerate}
\item<2-> Form \( BG \), study maps \( I:BG\to B\rr \) (invariant polynomials such as det, trace)
\item<3-> Fix a principal bundle \( P:M\to BG \)
\item<4-> Choose a connection \( A \) on \( P \), compute curvature \( F_A \)
\item<5-> Compose \( I\circ F_A \) (cohomology class)
\item<6-> \( I \) guarantees independence of class on connection
\end{enumerate}
\end{frame}

\begin{frame}
\begin{itemize}
\item What is \( BG\to B\rr \) in HoTT?
\item What are bundles, connections, and curvature in HoTT?
\item What is cohomology; what are exact forms?
\end{itemize}
\end{frame}

\section{Motivation}

\begin{frame}{Questions}
Where is the Bianchi identity?
What is ``the space of connections?''
\end{frame}

\section{Discrete manifolds}

\begin{frame}{Simplicial complexes}
% https://q.uiver.app/#q=WzAsNCxbMCwxLCJNXzA9XFxtYXRoYmJ7TX1fMCJdLFsxLDEsIlxcbWF0aGJie019XzEiXSxbMSwwLCJNXzEiXSxbMCwwLCJNXzFcXHRpbWVzIFxcYmRzaW1wbGV4bnsxfSJdLFswLDFdLFszLDAsIlxcbWF0aGJie0F9XzAiLDJdLFszLDIsIlxcbWF0aHJte3ByfV8xIl0sWzIsMSwiKl97XFxtYXRoYmJ7TX1fMX0iXSxbMSwzLCIiLDIseyJvZmZzZXQiOjMsInN0eWxlIjp7Im5hbWUiOiJjb3JuZXItaW52ZXJzZSJ9fV0sWzAsMiwiaF8xIiwyLHsic2hvcnRlbiI6eyJzb3VyY2UiOjQwLCJ0YXJnZXQiOjQwfSwibGV2ZWwiOjJ9XV0=
\begin{tikzcd}
  {M_1\times \bdsimplexn{1}} & {M_1} \\
  {M_0=\mathbb{M}_0} & {\mathbb{M}_1}
  \arrow["{\mathrm{pr}_1}", from=1-1, to=1-2]
  \arrow["{\mathbb{A}_0}"', from=1-1, to=2-1]
  \arrow["{*_{\mathbb{M}_1}}", from=1-2, to=2-2]
  \arrow["{h_1}"', shorten <=14pt, shorten >=14pt, Rightarrow, from=2-1, to=1-2]
  \arrow[from=2-1, to=2-2]
  \arrow["\ulcorner"{pos=0, rotate=180}, shift left=1, draw=none, from=2-2, to=1-1]
\end{tikzcd}

\end{frame}

\begin{frame}{Simplicial complexes}
% https://q.uiver.app/#q=WzAsNCxbMCwxLCJQKDIpXzAiXSxbMSwxLCJcXHBhcnRpYWwgXFxEZWx0YV4yIl0sWzEsMCwiXFxwYXJ0aWFsIFAoMikiXSxbMCwwLCJcXHBhcnRpYWwgUCgyKVxcdGltZXMgXFxwYXJ0aWFsIFxcRGVsdGFeMSJdLFswLDFdLFszLDAsIlxcbWF0aGJie0F9XzAiLDJdLFszLDIsIlxcbWF0aHJte3ByfV8xIl0sWzIsMSwiKl97XFxwYXJ0aWFsXFxEZWx0YV4yfSJdLFsxLDMsIiIsMix7InN0eWxlIjp7Im5hbWUiOiJjb3JuZXItaW52ZXJzZSJ9fV0sWzAsMiwiaF8xIiwyLHsic2hvcnRlbiI6eyJzb3VyY2UiOjQwLCJ0YXJnZXQiOjQwfSwibGV2ZWwiOjJ9XV0=
\begin{tikzcd}
  {\partial \Delta(2)\times \bdsimplexn{1}} & {\partial \Delta(2)} \\
  {\Delta(2)_0} & {\bdsimplex{2}}
  \arrow["{\mathrm{pr}_1}", from=1-1, to=1-2]
  \arrow["{\mathbb{A}_0}"', from=1-1, to=2-1]
  \arrow["{*_{\bdsimplexn{2}}}", from=1-2, to=2-2]
  \arrow["{h_1}"', shorten <=11pt, shorten >=11pt, Rightarrow, from=2-1, to=1-2]
  \arrow[from=2-1, to=2-2]
  \arrow["\ulcorner"{anchor=center, pos=0.125, rotate=180}, draw=none, from=2-2, to=1-1]
\end{tikzcd}

\end{frame}

\begin{frame}{Simplicial complexes}
% https://q.uiver.app/#q=WzAsNyxbMCwyLCJNXzFcXHRpbWVzIFxcYmRzaW1wbGV4bnsxfSJdLFswLDEsIlxcbWF0aGJie019XzAiXSxbMSwxLCJcXG1hdGhiYntNfV8xIl0sWzEsMiwiTV8xIl0sWzEsMCwiTV8yXFx0aW1lcyBcXGJkc2ltcGxleG57Mn0iXSxbMiwwLCJNXzIiXSxbMiwxLCJcXG1hdGhiYntNfV8yIl0sWzAsMSwiXFxtYXRoYmJ7QX1fMCJdLFsxLDJdLFswLDMsIlxcbWF0aHJte3ByfV8xIiwyXSxbMywyLCIqX3tcXG1hdGhiYntNfV8xfSIsMl0sWzIsMCwiIiwxLHsic3R5bGUiOnsibmFtZSI6ImNvcm5lci1pbnZlcnNlIn19XSxbMiw2XSxbNCwyLCJcXG1hdGhiYntBfV8xIiwyXSxbNSw2LCIqX3tcXG1hdGhiYntNfV8yfSJdLFs0LDUsIlxcbWF0aHJte3ByfV8xIl0sWzYsNCwiIiwxLHsic3R5bGUiOnsibmFtZSI6ImNvcm5lci1pbnZlcnNlIn19XSxbMiw1LCJoXzIiLDAseyJzaG9ydGVuIjp7InNvdXJjZSI6NDAsInRhcmdldCI6NDB9LCJsZXZlbCI6Mn1dLFsxLDMsImhfMSIsMix7InNob3J0ZW4iOnsic291cmNlIjo0MCwidGFyZ2V0Ijo0MH0sImxldmVsIjoyfV1d
\begin{tikzcd}
  & {M_2\times \bdsimplexn{2}} & {M_2} \\
  {\mathbb{M}_0} & {\mathbb{M}_1} & {\mathbb{M}_2} \\
  {M_1\times \bdsimplexn{1}} & {M_1}
  \arrow["{\mathrm{pr}_1}", from=1-2, to=1-3]
  \arrow["{\mathbb{A}_1}"', from=1-2, to=2-2]
  \arrow["{*_{\mathbb{M}_2}}", from=1-3, to=2-3]
  \arrow[from=2-1, to=2-2]
  \arrow["{h_1}"', shorten <=27pt, shorten >=27pt, Rightarrow, from=2-1, to=3-2]
  \arrow["{h_2}", shorten <=17pt, shorten >=17pt, Rightarrow, from=2-2, to=1-3]
  \arrow[from=2-2, to=2-3]
  \arrow["\ulcorner"{anchor=center, pos=0.125, rotate=-135}, draw=none, from=2-2, to=3-1]
  \arrow["\ulcorner"{anchor=center, pos=0.125, rotate=180}, draw=none, from=2-3, to=1-2]
  \arrow["{\mathbb{A}_0}", from=3-1, to=2-1]
  \arrow["{\mathrm{pr}_1}"', from=3-1, to=3-2]
  \arrow["{*_{\mathbb{M}_1}}"', from=3-2, to=2-2]
\end{tikzcd}

\end{frame}

\section{Codomain}

\begin{frame}{Codomain}
\( \EMzo \), \( \Kzt \)
\end{frame}

\section{Results}
\begin{frame}{Classical proof}
\begin{columns}
\column{0.5\textwidth}
\begin{figure}
\includegraphics[width=0.9\textwidth]{figs/needham_triangle.pdf}
\caption{Tristan Needham, \emph{Visual Differential Geometry and Forms}}
\end{figure}
\column{0.5\textwidth}
Every aspect was informed by the classical proof, which is already discrete-flavored. The ``\( \angle Fw_{||} \)'' looked a lot like a pathover.
\end{columns}
\end{frame}


\end{document}
